%% kapa: exam=1 round=36 year=2025 period=2 form=A subject=accounting subject_ko=회계학 seq=9 number=49
%% source: Q-Net 감정평가사 기출문제 (cst00309) artlSeq=5231525 file=raw/1008/5231525/2025년 제36회 감정평가사 1차시험 2교시 문제지 원본.pdf
%% conversion: kordoc 4.24.51 -> tools/gen1.py (auto); review=proofread-sonnet; flags=none
\begin{question}{49}
무형자산에 관한 설명으로 옳은 것은?
\begin{choices}
  \item 내용연수가 비한정인 무형자산으로 최초 인식한 경우 그 이후에 비한정 내용연수를 유한 내용연수로 변경할 수 없다.
  \item 원가모형과 달리 무형자산에 재평가모형을 적용하는 경우에는 원가가 아닌 금액으로 무형자산을 최초로 인식하는 것을 허용한다.
  \item 계약상 권리 또는 기타 법적 권리로부터 발생하는 무형자산의 내용연수는 자산의 예상 사용기간에 따라 그러한 계약상 권리 또는 기타 법적 권리의 기간을 초과할 수 있다.
  \item 제조과정에서 사용된 무형자산의 상각은 일반적으로 당기손익으로 인식한다.
  \item 자산에서 발생하는 미래경제적효익이 기업에 유입될 가능성이 높고 자산의 원가를 신뢰성 있게 측정할 수 있는 경우에만 무형자산을 인식한다.
\end{choices}
\end{question}
\answer{5}
